\section{Aislamiento electrico}

Para comprobar el aislamiento el\'ectrico del transformador, se conect\'o una
resistencia \(R_m\) entre el generador y el bobinado de entrada del
transformador. Luego se observ\'o simult\'aneamente la tensi\'on sobre dicha
resistencia, \(V_{Rm}\), y la tensi\'on de salida del transformador, \(V_{out}\),
al modificar el offset de continua de la se\~nal de entrada.

El objetivo de esta medici\'on fue verificar que el transformador transmite la
componente alterna de la se\~nal por acoplamiento magn\'etico, pero no permite
la transferencia directa de una componente continua entre el primario y el
secundario.

El valor nominal y el valor medido de la resistencia utilizada fueron:

\begin{table}[H]
\centering
\begin{tabular}{|c|c|c|}
\hline
Elemento & Valor nominal & Valor medido \\
\hline
\(R_m\) & \(47\,\Omega\) & \(50\,\Omega\) \\
\hline
\end{tabular}
\caption{Valor nominal y medido de la resistencia \(R_m\).}
\end{table}

\subsection{Medicion con offset de \(+5\,\mathrm{V}\)}

\begin{table}[H]
\centering
\begin{tabular}{|c|c|}
\hline
Magnitud & Valor medido \\
\hline
\(V_{Rm,pp}\) & \(380\,\mathrm{mV}\) \\
\hline
\(V_{Rm,max}\) & \(240\,\mathrm{mV}\) \\
\hline
\(V_{out,pp}\) & \(170\,\mathrm{mV}\) \\
\hline
\(V_{out,max}\) & \(84\,\mathrm{mV}\) \\
\hline
\end{tabular}
\caption{Mediciones obtenidas para un offset de \(+5\,\mathrm{V}\).}
\end{table}

\subsection{Medicion con offset de \(0\,\mathrm{V}\)}

\begin{table}[H]
\centering
\begin{tabular}{|c|c|}
\hline
Magnitud & Valor medido \\
\hline
\(V_{Rm,pp}\) & \(380\,\mathrm{mV}\) \\
\hline
\(V_{Rm,max}\) & \(169\,\mathrm{mV}\) \\
\hline
\(V_{out,pp}\) & \(170\,\mathrm{mV}\) \\
\hline
\(V_{out,max}\) & \(84\,\mathrm{mV}\) \\
\hline
\end{tabular}
\caption{Mediciones obtenidas para un offset de \(0\,\mathrm{V}\).}
\end{table}

\subsection{Medicion con offset de \(-5\,\mathrm{V}\)}

\begin{table}[H]
\centering
\begin{tabular}{|c|c|}
\hline
Magnitud & Valor medido \\
\hline
\(V_{Rm,pp}\) & \(380\,\mathrm{mV}\) \\
\hline
\(V_{Rm,max}\) & \(105\,\mathrm{mV}\) \\
\hline
\(V_{out,pp}\) & \(170\,\mathrm{mV}\) \\
\hline
\(V_{out,max}\) & \(84\,\mathrm{mV}\) \\
\hline
\end{tabular}
\caption{Mediciones obtenidas para un offset de \(-5\,\mathrm{V}\).}
\end{table}

Al comparar los tres casos, se observa que el valor de \(V_{Rm,max}\) var\'ia
al modificar el offset de continua aplicado por el generador. Sin embargo, la
tensi\'on de salida del transformador permanece pr\'acticamente constante:
\(V_{out,pp}=170\,\mathrm{mV}\) y \(V_{out,max}=84\,\mathrm{mV}\) para los tres
valores de offset ensayados.

Adem\'as, el valor de \(V_{out,max}\) resulta aproximadamente igual a la mitad
de \(V_{out,pp}\), ya que:

\[
\frac{V_{out,pp}}{2}=\frac{170\,\mathrm{mV}}{2}=85\,\mathrm{mV}
\]

Este resultado indica que la se\~nal de salida se mantiene centrada alrededor
de cero, sin arrastrar el desplazamiento de continua presente en la entrada.
Por lo tanto, el transformador bloquea la componente continua y solamente
transfiere la componente alterna de la se\~nal.

